\documentclass[12pt]{article}
\usepackage[a4paper,margin=1in]{geometry}
\usepackage[utf8]{inputenc}
\usepackage[T1]{fontenc}
\usepackage{amsmath,amssymb}
\usepackage{textcomp}
\usepackage{parskip}
\begin{document}
\begin{center}
{\Large \textbf{2006 STEP III}}
\end{center}

\subsection*{Page 1}
\begin{flushleft}
STEP III, 2006 2\par
\bigskip\textbf{Section A: Pure Mathematics}\par
\bigskip\textbf{Question 1.} Sketch the curve with cartesian equation\par
\[2x(x^2 5)\]
\[y = -\]
\[x^2 4\]
-\par
and give the equations of the asymptotes and of the tangent to the curve at the origin.\par
Hence determine the number of real roots of the following equations:\par
\textbf{(i)} 3x(x2 5) = (x2 4)(x + 3) ;\par
- -\par
\textbf{(ii)} 4x(x2 5) = (x2 4)(5x 2) ;\par
- - -\par
\textbf{(iii)} 4x2(x2 5)2 = (x2 4)2(x2 + 1) .\par
- -\par
2 Let\par
\[1 2 \pi cos2 θ 1 2 \pi sec2 θ\]
\[I = dθ and J = dθ\]
\[Z- 1 \pi 1 sin θ sin 2α Z- 1 \pi 1 + tan2 θ cos2 2α\]
2 - 2\par
\[where 0 < α < \pi .\]
\[1 2 \pi cos2 θ 2 1 \pi 2\]
\textbf{(i)} Show that I = dθ and hence that 2I = dθ .\par
\[Z- 1 \pi 1 + sin θ sin 2α Z- 1 \pi 1 + tan2 θ cos2 2α\]
2 2\par
\textbf{(ii)} Find J .\par
\textbf{(iii)} By considering I sin 2 2α + J cos2 2α, or otherwise, show that I = 1 \pi sec2 α.\par
1 1\par
\textbf{(iv)} Evaluate I in the case \pi < α < \pi.\par
4 2\par
\end{flushleft}
\subsection*{Page 2}
\begin{flushleft}
STEP III, 2006 3\par
\bigskip\textbf{Question 3.} (i) Let\par
\infty \infty\par
n n\par
\[tan x = anx and cot x = + bnx\]
x\par
\[nX=0 nX=0\]
\[for 0 < x <\]
\[\pi . Explain why an = 0 for even n.\]
Prove the identity\par
cot x tan x 2 cot 2x\par
- \equiv\par
and show that\par
n+1\par
\[an = (1 2 )bn .\]
-\par
\infty\par
n 1\par
\textbf{(ii)} Let cosec x =\par
x\par
\[+ cnx for 0 < x <\]
\[\pi . By considering cot x + tan x, or otherwise,\]
\[Xn=0\]
show that\par
-n\par
\[cn = (2 1)bn .\]
-\par
\textbf{(iii)} Show that\par
\infty 2 \infty 2\par
n 2 n\par
\[1 + x bnx + x = 1 + x cnx .\]
\[\lambda n=0 \mu \lambda n=0 \mu\]
P P\par
\[Deduce from this and the previous results that a^1 = 1, and find a^3.\]
\bigskip\textbf{Question 4.} The function f satisfies the identity\par
f(x) + f(y) f(x + y) ( )\par
\equiv ∗\par
′′\par
\[for all x and y. Show that 2f(x) f(2x) and deduce that f (0) = 0. By considering the\]
\equiv\par
Maclaurin series for f(x), find the most general function that satisfies ( ).\par
∗\par
[Do not consider issues of existence or convergence of Maclaurin series in this question.]\par
\textbf{(i)} By considering the function G, defined by ln g(x) = G(x), find the most general\par
function that, for all x and y, satisfies the identity\par
g(x)g(y) g(x + y) .\par
\equiv\par
u\par
\textbf{(ii)} By considering the function H, defined by h(e ) = H(u), find the most general function\par
that satisfies, for all positive x and y, the identity\par
h(x) + h(y) h(xy) .\par
\equiv\par
\textbf{(iii)} Find the most general function t that, for all x and y, satisfies the identity\par
t(x) + t(y) t(z) ,\par
\equiv\par
x + y\par
\[where z = .\]
1 xy\par
-\par
\end{flushleft}
\subsection*{Page 3}
\begin{flushleft}
STEP III, 2006 4\par
\bigskip\textbf{Question 5.} Show that the distinct complex numbers α, β and γ represent the vertices of an equilateral\par
triangle (in clockwise or anti-clockwise order) if and only if\par
2 2 2\par
\[α + β + γ βγ γα αβ = 0 .\]
- - -\par
Show that the roots of the equation\par
3 2\par
\[z + az + bz + c = 0 ( )\]
∗\par
\[represent the vertices of an equilateral triangle if and only if a^2 = 3b.\]
\[Under the transformation z = pw + q, where p and q are given complex numbers with p = 0,\]
the equation ( ) becomes\par
∗ 3 2\par
\[w + Aw + Bw + C = 0 . ( )\]
∗∗\par
Show that if the roots of equation ( ) represent the vertices of an equilateral triangle, then\par
∗\par
the roots of equation ( ) also represent the vertices of an equilateral triangle.\par
∗∗\par
\bigskip\textbf{Question 6.} Show that in polar coordinates the gradient of any curve at the point (r, θ) is\par
dr\par
tan θ + r\par
dθ\par
.\par
dr\par
r tan θ\par
dθ -\par
A mirror is designed so that if an incident ray of light is parallel to a fixed line L the reflected\par
ray passes through a fixed point O on L. Prove that the mirror intersects any plane containing\par
L in a parabola. You should assume that the angle between the incident ray and the normal\par
to the mirror is the same as the angle between the reflected ray and the normal.\par
L O\par
\end{flushleft}
\subsection*{Page 4}
\begin{flushleft}
STEP III, 2006 5\par
\bigskip\textbf{Question 7.} (i) Solve the equation u2 + 2u sinh x 1 = 0 giving u in terms of x.\par
-\par
Find the solution of the differential equation\par
dy dy\par
\[+ 2 sinh x 1 = 0\]
\lambda dx \mu dx -\par
dy\par
\[that satisfies y = 0 and > 0 at x = 0.\]
dx\par
\textbf{(ii)} Find the solution, not identically zero, of the differential equation\par
dy dy\par
\[sinh y + 2 sinh y = 0\]
\lambda dx \mu dx -\par
\[that satisfies y = 0 at x = 0, expressing your solution in the form cosh y = f(x). Show\]
\[that the asymptotes to the solution curve are y = ( x + ln 4).\]
± -\par
8 is an operation that takes polynomials in x to polynomials in x; that is, given any polynomial\par
△\par
h(x), there is a polynomial called h(x) which is obtained from h(x) using the rules that define\par
△\par
. These rules are as follows:\par
△\par
\textbf{(i)} x = 1 ;\par
△\par
\textbf{(ii)} f(x) + g(x) = f(x) + g(x) for any polynomials f(x) and g(x);\par
△ △ △\par
\textbf{(iii)} λf(x) = λ f(x) for any constant λ and any polynomial f(x);\par
△ △\par
\textbf{(iv)} f(x)g(x) = f(x) g(x) + g(x) f(x) for any polynomials f(x) and g(x).\par
△ △ △\par
Using these rules show that, if f(x) is a polynomial of degree zero (that is, a constant), then\par
\[f(x) = 0. Calculate x^2 and x^3.\]
△ △ △\par
dh(x)\par
Prove that h(x) for any polynomial h(x). You should make it clear whenever you\par
△ \equiv dx\par
use one of the above rules in your proof.\par
\end{flushleft}
\subsection*{Page 5}
\begin{flushleft}
STEP III, 2006 6\par
\bigskip\textbf{Section B: Mechanics}\par
\bigskip\textbf{Question 9.} A long, light, inextensible string passes through a small, smooth ring fixed at the point O.\par
One end of the string is attached to a particle P of mass m which hangs freely below O. The\par
other end is attached to a bead, B, also of mass m, which is threaded on a smooth rigid wire\par
fixed in the same vertical plane as O. The distance OB is r, the distance OH is h and the\par
height of the bead above the horizontal plane through O is y, as shown in the diagram.\par
H\par
B\par
h\par
r\par
y\par
θ\par
O\par
P\par
The shape of the wire is such that the system can be in static equilibrium for all positions of\par
the bead. By considering potential energy, show that the equation of the wire is y + r = 2h.\par
The bead is initially at H. It is then projected along the wire with initial speed V . Show\par
that, in the subsequent motion,\par
hr˙\par
\[θ˙ =\]
\[- r\sqrt\{rh\} h^2\]
-\par
\[where θ is given by θ = arcsin(y/r).\]
r h\par
Hence show that the speed of the particle P is V - .\par
2r h\par
\gamma \delta\par
-\par
[Note that arcsin θ is another notation for sin\par
-1\par
θ.]\par
\end{flushleft}
\subsection*{Page 6}
\begin{flushleft}
STEP III, 2006 7\par
\bigskip\textbf{Question 10.} A disc rotates freely in a horizontal plane about a vertical axis through its centre. The moment\par
of inertia of the disc about this axis is mk2 (where k > 0). Along one diameter is a smooth\par
narrow groove in which a particle of mass m slides freely. At time t = 0 , the disc is rotating\par
with angular speed Ω, and the particle is a distance a from the axis and is moving with speed V\par
\[along the groove, towards the axis, where k2V 2 = Ω2a2(k^2 + a^2) .\]
Show that, at a later time t, while the particle is still moving towards the axis, the angular\par
speed ω of the disc and the distance r of the particle from the axis are related by\par
\[Ω(k^2 + a^2) dr 2 Ω2r2(k^2 + a^2)2\]
\[ω = and = .\]
\[k^2 + r^2 \lambda dt \mu k^2(k^2 + r^2)\]
Deduce that\par
dr 2 2 1\par
\[k = r(k + r ) 2 ,\]
dθ -\par
where θ is the angle through which the disc has turned by time t.\par
\[By making the substitution u = k/r, or otherwise, show that r sinh(θ + α) = k, where\]
\[sinh α = k/a . Deduce that the particle never reaches the axis.\]
\bigskip\textbf{Question 11.} A lift of mass M and its counterweight of mass M are connected by a light inextensible cable\par
which passes over a fixed frictionless pulley. The lift is constrained to move vertically between\par
smooth guides. The distance between the floor and the ceiling of the lift is h. Initially, the\par
lift is at rest, and the distance between the top of the lift and the pulley is greater than h. A\par
small tile of mass m becomes detached from the ceiling of the lift and falls to the floor of the\par
lift. Show that the speed of the tile just before the impact is\par
(2M m)gh\par
- .\par
r M\par
The coefficient of restitution between the tile and the floor of the lift is e. Given that the\par
magnitude of the impulsive force on the lift due to tension in the cable is equal to the magnitude\par
of the impulsive force on the counterweight due to tension in the cable, show that the loss of\par
\[energy of the system due to the impact is mgh(1 e^2). Comment on this result.\]
-\par
\end{flushleft}
\subsection*{Page 7}
\begin{flushleft}
STEP III, 2006 8\par
\bigskip\textbf{Section C: Probability and Statistics}\par
12 Fifty times a year, 1024 tourists disembark from a cruise liner at a port. From there they must\par
travel to the city centre either by bus or by taxi. Tourists are equally likely to be directed to\par
the bus station or to the taxi rank. Each bus of the bus company holds 32 passengers, and\par
the company currently runs 15 buses. The company makes a profit of \pounds 1 for each passenger\par
carried. It carries as many passengers as it can, with any excess being (eventually) transported\par
by taxi. Show that the largest annual licence fee, in pounds, that the company should consider\par
paying to be allowed to run an extra bus is approximately\par
800\par
-2\par
1600Φ(2) 1 e ,\par
\[- \sqrt\{2\}\pi -\]
x\par
\[1 - 1 t^2\]
\[where Φ(x) = e 2 dt .\]
\[\sqrt\{2\}\pi Z-\infty\]
[You should not consider continuity corrections.]\par
\bigskip\textbf{Question 13.} Two points are chosen independently at random on the perimeter (including the diameter) of\par
a semicircle of unit radius. The area of the triangle whose vertices are these two points and\par
the midpoint of the diameter is denoted by the random variable A. Show that the expected\par
\[value of A is (2 + \pi)\]
-1.\par
\bigskip\textbf{Question 14.} For any random variables X1 and X2, state the relationship between E(aX1 + bX2) and E(X1)\par
and E(X2), where a and b are constants. If X1 and X2 are independent, state the relationship\par
\[between E(X1X2) and E(X^1) and E(X^2).\]
An industrial process produces rectangular plates. The length and the breadth of the plates\par
\[are modelled by independent random variables X^1 and X^2 with non-zero means µ1 and µ2 and\]
non-zero standard deviations σ1 and σ2, respectively. Using the results in the paragraph above,\par
\[and without quoting a formula for Var (aX1 + bX2), find the means and standard deviations\]
of the perimeter P and area A of the plates. Show that P and A are not independent.\par
\[The random variable Z is defined by Z = P αA, where α is a constant. Show that Z and A\]
-\par
are not independent if\par
2(µ σ2 + µ σ2)\par
1 2 2 1\par
\[α = .\]
6 µ2σ2 + µ2σ2 + σ2σ2\par
1 2 2 1 1 2\par
Given that X1 and X2 can each take values 1 and 3 only, and that they each take these values\par
with probability , show that Z and A are not independent for any value of α.\par
\end{flushleft}
\end{document}
